%%\documentstyle[twocolumn,jsaiac]{jarticle}
%%\documentstyle[twocolumn,jsaiac]{j-article}
\documentclass[twocolumn]{jarticle}

\usepackage{jsaiac}

\usepackage{color}
\usepackage{url}
\usepackage{booktabs}
\usepackage{framed}
\usepackage{multicol}
\usepackage{multirow}

%%
\title{
\jtitle{LA-Bench 2025：実験指示から実行可能手順を生成するための\\データセット}
\etitle{LA-Bench 2025: A Dataset for Generating Executable Experimental Procedures \\from Experimental Instructions}
}
%%英文は以下を使用
%\title{Style file for manuscripts of JSAI 20XX}

\jaddress{山田 涼太，Science Aid株式会社，東京都中央区日本橋本町2丁目3番11号，ryota.yamada@science-aid.com}

\author{%
\jname{加藤 祥太\first}
\ename{Shota Kato}
\and
\jname{西田 理彦\second}
\ename{Toshihiko Nishida}
\and
\jname{酒井 雄介\third}
\ename{Yusuke Sakai}
\and
\jname{尾崎 遼\fourth}
\ename{Haruka Ozaki}
\and
\jname{杉山 亜矢斗\second}
\ename{Ayato Sugiyama}
\and
\jname{山田 涼太\fifth}
\ename{Ryota Yamada}
%Given-name Surname\third{}%%英文は左を使用
}

\affiliate{
\jname{\first{}京都大学}
\ename{Kyoto University}
\and
\jname{\second{}ラボラトリーオートメーション協会}
\ename{Laboratory Automation Suppliers' Association}
\and
\jname{\third{}東京大学}
\ename{The University of Tokyo}
\and
\jname{\fourth{}理化学研究所}
\ename{RIKEN}
\and
\jname{\fifth{}Science Aid株式会社}
\ename{Science Aid}
}

%%
%\Vol{28}        %% <-- 28th（変更しないでください）
%\session{0A0-00}%% <-- 講演ID（必須)

\begin{abstract}
実験指示から実行可能な実験手順を生成する能力を評価するため，ベンチマークデータセットLA-Bench 2025を構築し，大規模言語モデル（LLM）を評価器として用いるLLM-as-a-judgeの評価設計を整理した．評価は共通採点基準（5点）と個別採点基準（5点）の合計10点満点とし，設計思想の異なる3種類の評価プロンプトと3種類のLLM評価器の組み合わせで計450回の評価を行った．その結果，得点傾向はプロンプトと評価器の組み合わせに依存し，最高得点を与えるプロンプトが評価器の種類によって入れ替わることを確認した．以上より，実験手順生成タスクの自動評価では単一の評価設計に依存せず，複数の評価設計を併用する必要があることが示唆された．
\end{abstract}

%\setcounter{page}{1}
\def\Style{``jsaiac.sty''}
\def\BibTeX{{\rm B\kern-.05em{\sc i\kern-.025em b}\kern-.08em%
 T\kern-.1667em\lower.7ex\hbox{E}\kern-.125emX}}
\def\JBibTeX{\leavevmode\lower .6ex\hbox{J}\kern-0.15em\BibTeX}
\def\LaTeXe{\LaTeX\kern.15em2$_{\textstyle\varepsilon}$}

\begin{document}
\maketitle

\section{はじめに}
自動実験装置と機械学習を組み合わせ，実験の計画から実行・解析までを閉ループで回す自律実験は，材料・化学・生命科学の各分野で研究が進んでいる~\cite{hase2019}．この文脈では，「どの実験を次に実行するか」という探索・最適化戦略の性能を評価するベンチマークの整理が進んでいる．例えば，自律駆動ラボ（self-driving laboratory; SDL）のベンチマークに関するレビューでは，探索戦略をランダム探索のような参照手法と比較し，加速係数（acceleration factor; AF）および改善係数（enhancement factor; EF）によって定量化する枠組みが示されている~\cite{benchmarking_sdl_2026}．
一方で，自律実験の実装においては，探索よりも前段の工程として，目的や条件の記述といった実験指示を，自動実験装置が実行可能な粒度の実験手順へ変換する工程が存在する~\cite{shi2024expert}．この工程では，物品制約，操作順序，数量整合，安全上の注意，参考手順との整合性など，多面的な制約を同時に満たす必要がある．本研究では，この実験指示から実行可能手順を生成する工程を対象とする．

実験手順生成を対象とした評価は，人手評価が高コストであることに加え，合理的な別解が存在するため，単純な文字列一致や構造一致では妥当な評価が困難である．
ウェット実験プロトコルを対象とした情報抽出研究として，WLPデータセット~\cite{kulkarni2018wlp}やWNUT-2020タスク~\cite{tabassum2020wnut}，材料合成手順を構造化するPcMSP~\cite{yang2022pcmsp}がある．しかし，これらは既存手順の抽出・構造化を主目的とするものであり，実験指示から実行可能手順を生成する能力の評価とは異なる．
近年は大規模言語モデル（large language model; LLM）を評価器として用いるLLM-as-a-judgeが提案されているが~\cite{zheng2023mtbench,liu2023geval}，実験手順のように安全性と実行可能性が重要な領域において，評価観点と評価プロンプトを明示的に設計し，評価対象のレイヤーを整理したベンチマークは存在しない．

このような背景を踏まえ，本研究では，実験指示から実行可能な実験手順を生成する能力を評価するためのデータセット\textbf{LA-Bench 2025}を構築した．
LA-Bench 2025は，実験目的や条件を記述した自然言語の実験指示，自動実験装置で実行可能な粒度の実験手順，および採点基準から構成される．
本稿では，LA-Bench 2025のデータ設計および評価設計を整理し，LLM-as-a-judgeに基づく評価プロンプトの構成を明示する．さらに，評価プロンプトと評価器（LLM）の選択が得点に与える影響を分析し，実験手順生成タスクにおける評価設計の重要性を述べる．

\section{方法}
\subsection{タスク定義とデータ構造}
LA-Bench 2025は，実験目的や条件を記述した自然言語の実験指示から，自動実験装置で実行可能な粒度の実験手順を生成するタスクである．
各問題は構造化された入力情報を持ち，それに対応する番号付き手順列を出力とする．
入出力の構成を表\ref{tab:io_format}に示す．
\begin{table}[t]
    \centering
    \caption{LA-Bench 2025における入出力情報}
    \label{tab:io_format}
    \begin{tabular}{lp{6.5cm}}
    \toprule
    区分 & 内容 \\
    \midrule
    入力 & 1. 実験指示（実験目的と要求事項）\\
    & 2. 必ず使用する物品の一覧 \\
    & 3. 元プロトコルの手順（参考となる実験手順）\\
    & 4. 期待される最終状態（最終的に得られるサンプルや重要な中間生成物とその状態）\\
    & 5. 参考文献 \\
    \midrule
    出力 & 番号付きの実験手順．各手順は具体的な数量・条件・操作対象を明示し，学部レベルのウェット実験の訓練を受けた者が誤解なく実行可能な粒度とする．手順数は50以下，各手順は10文以下とする．\\
    \bottomrule
    \end{tabular}
\end{table}
入力の主要な要素は評価基準と対応している．必ず使用する物品の一覧は手順中に適切に使用されていることが求められ，元プロトコルの手順は論理構造の整合性を評価する基準となる．また，期待される最終状態は，生成手順が実験として完遂可能であるかを判断するための参照情報として機能する．
出力手順は単なる構造化ではなく，暗黙知の補完，数値パラメータ（体積・濃度・温度・時間など）の明示，操作対象や装置の特定を含むことを求める．
問題数は，例題セット5問，開発用テストセット10問，最終評価用テストセット10問の合計25問である．
各問題は，物品制約や論理構造（順序・分岐）を含む異なる操作特性を持つように，対象とする実験手順の専門家が作成した．
データと詳細仕様は公開リポジトリに示す\footnote{\url{https://github.com/lasa-or-jp/la-bench}}．

\subsection{評価設計}
評価は，全問題に共通する共通採点基準（5点満点）と，問題ごとに定義される個別採点基準（5点満点）の合計10点満点で行う．
いずれの採点基準も上記の専門家が設計した．

共通採点基準は，一般的な実行可能性を評価する．具体的には，(1) 実験指示に含まれるパラメータが手順に正確に反映されているか，(2) 必須物品が適切に使用されているか，(3) 元プロトコルの論理構造（順序・分岐）が保持されているか，(4) 手順全体として期待される最終状態を論理的に達成可能か，(5) 明示されていない暗黙部分が適切に補完されているか，を評価する．
一方，計算誤り，論理矛盾，不自然な記述が認められる場合には減点する．減点は各項目につき1回のみ適用する．また，入力手順の単純転記や過度な安全側への条件拡張など，実質的な思考を伴わない出力と判断される場合には，共通採点基準の合計点を最大2点に制限する．この制限は，過度の安全性による評価の歪みを抑制するためのものである．

個別採点基準は，各問題の出題意図に基づいて設定される．安全性，効率性，精度向上，コスト削減などの観点から具体的な達成条件を定義し，合計5点となるよう配点する．

\subsection{自動評価と評価プロンプト}\label{sec:judge_design}
自動評価にはLLM-as-a-judgeを用いる．共通採点基準および個別採点基準に基づき，各回答に対して共通得点（0--5点）と個別得点（0--5点）を算出し，その合計を最終得点（0--10点）とする．また，共通採点基準の判定理由，個別採点基準に該当した根拠，および特記事項を併せて出力させる．

評価プロンプトは，システムプロンプトと，入力・出力・個別採点基準を埋め込む共通テンプレートから構成される．
評価プロンプト設計に起因する偏りを低減するため，設計思想の異なる3種類のシステムプロンプト（厳格なチェックリスト型，自律推論型，対照デュアルパス型）を独立に用い，それぞれで評価を行う．

厳格なチェックリスト型プロンプトでは，各共通採点基準を二値（OK/NG）で独立に判定し，不確実な場合は無加点とする．判定根拠は出力手順からのみ引用させ，共通採点基準の判定理由を固定フォーマットで出力させることで，得点と根拠の対応を明示的に拘束する．また，入力手順の単純転記や過度な安全側拡張が認められる場合には，共通採点基準の合計点を最大2点とする過度の安全性制限を適用する．
厳格なチェックリスト型プロンプトのうち，共通採点基準，過度の安全性制限，および出力形式に関する部分を図\ref{fig:strict_prompt}に示す．
\begin{figure}[!t]
  \centering
  \begin{minipage}{0.8\columnwidth}
  \begin{framed}
      \begin{verbatim}
# 共通採点基準（最大5点）
1. 実験指示のパラメータを正確に反映
2. 必須物品をすべて具体的に使用
3. 元プロトコルの順序・分岐を保持
4. 最終状態を論理的に達成可能
5. 暗黙部分を適切に補完

# 過度の安全性
- 入力手順の単純転記
- 出典の無批判な転記
- 不必要な安全側拡張
→ 共通得点を最大2点に制限

# 出力形式
{
  "general_score": 0-5,
  "specific_score": 0-5,
  "final_score": 0-10,
  "general_reason": "...",
  "specific_matches": [...],
  "notes": "..."
}
\end{verbatim}
  \vspace{-\baselineskip}
  \end{framed}
  \end{minipage}
  \caption{厳格なチェックリスト型評価プロンプトの抜粋．共通採点基準，過度の安全性制限，および出力形式を示す．}
  \label{fig:strict_prompt}
\end{figure}

自律推論型プロンプトでは，評価手順を明示し，自律的に推論を進める設計とした．対照デュアルパス型プロンプトでは，失敗理由探索と加点根拠探索を分離し，両者の衝突解決規則を明示することで，曖昧な判定を抑制した．表\ref{tab:judge_prompt_comparison}に，各プロンプトの設計思想を示す．
\begin{table*}[t]
    \centering
    \caption{3種類の評価プロンプトの設計思想の比較}
    \label{tab:judge_prompt_comparison}
    \begin{tabular}{llp{8cm}}
        \toprule
        評価プロンプト & 設計思想 & 主な特徴 \\ \midrule
        厳格なチェックリスト型 & 二値判定による保守的評価 &
        各共通採点基準を独立に判定し，不確実な場合は無加点とする．
        過度の安全性制限を明示的に適用する． \\

        自律推論型 & 段階的推論による体系的評価 &
        入力抽出，出力解析，基準適用，集計の判断手順を明示して評価を行う．
        出力形式を制約し，冗長な出力を禁止する．\\

        対照デュアルパス型 & 失敗優先の対照評価 &
        失敗理由探索と加点根拠探索を分離する．
        衝突解決規則を明示し，曖昧な場合は無加点とする．\\ \bottomrule
    \end{tabular}
\end{table*}

\subsection{実験}
評価プロンプトの設計と評価器の選択が得点に与える影響を分析するため，開発用テストセット10問の正解手順を評価対象として自動評価を実施した．
評価プロンプトは\ref{sec:judge_design}節の3種類，評価器はGPT-5-nano，GPT-5-mini，GPT-5の3種類とした．
LLMのtemperatureパラメータは1.0とした．
各組み合わせ（$3 \times 3 = 9$ 通り）について各問題を5回ずつ評価し，その平均値を分析した．
実装の詳細は公開リポジトリを参照されたい\footnote{実装は\texttt{code/evaluation/run\_judge\_experiment.py}にある．}．


\section{結果と考察}
評価プロンプトと評価器の組み合わせごとの評価結果を表\ref{tab:results_matrix}に示す．
\begin{table*}[t]
    \centering
    \caption{評価プロンプトと評価器の組み合わせごとの評価結果（平均 $\pm$ 標準偏差）}
    \label{tab:results_matrix}
    \begin{tabular}{lcccc}
        \toprule
        \multirow{2}{*}{評価器} & \multicolumn{3}{c}{評価プロンプト} & \multirow{2}{*}{プロンプト平均}\\
        & チェックリスト型 & 自律推論型 & デュアルパス型 & \\
        \midrule
        GPT-5-nano & $8.10 \pm 2.62$ & $7.30 \pm 2.82$ & $7.48 \pm 2.61$ & $7.63 \pm 2.69$ \\
        GPT-5-mini & $7.54 \pm 2.23$ & $7.28 \pm 2.47$ & $7.18 \pm 2.66$ & $7.33 \pm 2.45$ \\
        GPT-5      & $8.00 \pm 1.59$ & $8.48 \pm 1.20$ & $7.96 \pm 1.78$ & $8.15 \pm 1.55$ \\
        \midrule
        プロンプト平均 & $7.88 \pm 2.19$ & $7.69 \pm 2.32$ & $7.54 \pm 2.39$ & $7.70 \pm 2.30$ \\
        \bottomrule
    \end{tabular}
\end{table*}
全体では平均7.70点，標準偏差2.30点であった．
問題ごとの平均得点の範囲は $[3.42, 9.33]$ であり，満点（10点）に近い問題は限定的であった．
共通採点基準の平均は 3.49 点，個別採点基準の平均は 4.22 点であり，個別採点基準の方が高い傾向が認められた．
この傾向は全9通りの評価器・プロンプトの組み合わせで一貫しており，正解手順は問題固有の要件を概ね満たす一方，共通採点基準が求める入力忠実性や論理的整合性の面では減点が生じた．
% --- 発表用補足データ: 共通(general) vs 個別(specific) 得点の組み合わせ別内訳（平均±標準偏差） ---
% GPT-5-nano x checklist  : general=3.98±1.36  specific=4.12±1.45  diff=+0.14
% GPT-5-nano x autonomous : general=3.52±1.54  specific=3.78±1.59  diff=+0.26
% GPT-5-nano x dual-pass  : general=3.62±1.44  specific=3.96±1.37  diff=+0.34
% GPT-5-mini x checklist  : general=3.30±1.57  specific=4.24±1.08  diff=+0.94
% GPT-5-mini x autonomous : general=3.02±1.56  specific=4.26±1.31  diff=+1.24
% GPT-5-mini x dual-pass  : general=2.94±1.81  specific=4.24±1.20  diff=+1.30
% GPT-5      x checklist  : general=3.62±1.09  specific=4.38±0.85  diff=+0.76
% GPT-5      x autonomous : general=3.94±0.74  specific=4.54±0.79  diff=+0.60
% GPT-5      x dual-pass  : general=3.50±1.36  specific=4.46±0.86  diff=+0.96
% 全9通りで specific > general が成立．
% GPT-5-mini で差が最大（+0.94〜+1.30），GPT-5-nano で差が最小（+0.14〜+0.34）．
% GPT-5 は general/specific ともに標準偏差が小さく，評価が安定している．
% ---

表\ref{tab:results_matrix}では，評価器とプロンプトの得点傾向が独立ではないことが読み取れる．GPT-5-nanoおよびGPT-5-miniではチェックリスト型が最高得点（それぞれ$8.10$, $7.54$）を示すのに対し，GPT-5では自律推論型が最高得点（$8.48$）を示している．
プロンプトの観点からは，チェックリスト型がプロンプト平均で最高（$7.88$）であったが，これは主にGPT-5-nanoおよびGPT-5-miniでの高得点に起因する．正解手順が各基準を明示的に満たしているため，二値判定による加点が安定して行われたと考えられる．対照デュアルパス型は全評価器で最低または最低に近い得点を示し（プロンプト平均$7.54$），失敗理由を優先的に探索する設計が減点根拠を見出しやすくした結果と解釈できる．

評価器（LLM）間の全体平均得点差は$0.81$点であった（GPT-5: $8.15$，GPT-5-nano: $7.63$，GPT-5-mini: $7.33$）．ただし，この差はプロンプトにより大きく異なる．評価器間得点差は，チェックリスト型で$0.56$点，自律推論型で$1.20$点，対照デュアルパス型で$0.78$点であり，自律推論型において最も大きかった．自律推論型は段階的推論を要求する設計であるため，推論能力の低いGPT-5-nanoとGPT-5-miniでは途中の判断を誤りやすく得点が低下する一方，推論能力の高いGPT-5は設計意図を適切に遂行でき高得点を示したと考えられる．チェックリスト型は二値判定の単純なタスクであるため，評価器間差が最も小さかった．ただし，この解釈の検証には専門家評価との対照が必要であり，今後の課題である．

以上の結果は，プロンプト設計と評価器選択が相互に影響し，評価結果に直接反映されることを示している．
単一のプロンプトや評価器に依存した評価では，特定の組み合わせに固有の偏りを排除できないため，複数の組み合わせによる評価が必要である．
共通採点基準における過度の安全性制限は，不必要に条件を安全側へ拡張する傾向を抑制するために導入した．
この設計により，入力忠実性と実行可能性のバランスを評価できるようになった．一方で，LLM-as-a-judgeは自然言語的整合性を評価する能力に優れるが，数量整合や形式的制約の検証においては限界を持つ．今後は，形式的検証手法やルールベースの整合性チェックとLLM評価を統合する枠組みが有望である．
また，評価指標が最適化対象となると，実験手順生成手法がその指標に特化する可能性があるため，評価設計はデータセットの進化とともに更新されることが望ましい．

\section{おわりに}
本稿では，実験指示から実行可能な実験手順を生成するタスクに対して，LA-Bench 2025のデータ設計および評価設計を整理し，LLM-as-a-judgeに基づく評価プロンプトの構成を明示した．また，評価プロンプトの設計および評価器の選択により得点が変動し得ることを示し，実験手順生成タスクにおける評価設計の重要性を述べた．

本データセットはコンペティション\footnote{\url{https://lasa-or-jp.github.io/la-bench/}}においても使用し，複数のグループから提出された生成手順に対して，自動評価および一部の提出に対する専門家による人手評価を実施した．
しかし，本稿では主としてデータと自動評価設計に焦点を当てたため，提出物全体の傾向分析や，自動評価結果と人手評価結果の差異分析，別解の扱いを含む評価の妥当性検証は今後の課題である．

\section*{謝辞}
LA-Bench 2025に参加し，実験手順生成手法を提出いただいた各チームの皆様に感謝する．
また，本取り組みは一般社団法人ラボラトリーオートメーション協会（Laboratory Automation Suppliers' Association; LASA）および人工知能学会のコンペティション開催支援制度の支援を受けた．さらに，スポンサーとして株式会社ディープコア，合同会社embrio，エピストラ株式会社，HERO Impact Capital，株式会社Inner Resource，株式会社システム計画研究所，株式会社リンカーズOI研究所，Polaris.AI株式会社，プライマルキャピタルの各団体にご支援をいただいた．ここに深く感謝の意を表する．

本研究の一部は，JST さきがけ JPMJPR25T1 の支援を受けたものである．

\bibliographystyle{jsai}
\bibliography{reference}

\end{document}
